% ==================================================================
% id:       cond-006
% topic:    הסתברות מותנית ואי-תלות
% subtopic: הסתברות מותנית
% solution: full
% ==================================================================
מטילים שתי קוביות הוגנות, אחת צהובה ואחת אדומה.
\begin{enumerate}
\item מה מרחב המדגם?
\begin{solbox}
זוגות סדורים (התוצאה בצהובה, התוצאה באדומה): $\Omega=\set{1,\dots,6}\times\set{1,\dots,6}$, ו-$|\Omega|=36$. הקוביות הוגנות, ולכן המרחב אחיד.
\end{solbox}
\item מה ההסתברות שבקובייה האדומה יצא 3?
\begin{solbox}
יש 6 תוצאות שבהן באדומה יצא 3, ולכן ההסתברות היא $\frac6{36}=\ans{\frac16}$.
\end{solbox}
\item מה ההסתברות שבשתי הקוביות יצא 3?
\begin{solbox}
רק התוצאה $(3,3)$, ולכן ההסתברות היא $\ans{\frac1{36}}$.
\end{solbox}
\item ידוע שבקובייה האדומה יצא 3. מה ההסתברות שבשתי הקוביות יצא 3?
\begin{solbox}
נסמן $A$ = „בשתי הקוביות 3”, $B$ = „באדומה 3”. כאן $A\subseteq B$, ולכן $A\cap B=A$, ולפי סעיפים (ב) ו-(ג):
\[
\PP(A\mid B)=\frac{\PP(A)}{\PP(B)}=\frac{1/36}{1/6}=\ans{\frac16}.
\]
זה הגיוני: אם באדומה כבר יצא 3, נשאר רק שבצהובה ייצא 3.
\end{solbox}
\item ידוע ששתי התוצאות שוות. מה ההסתברות ששתיהן זוגיות?
\begin{solbox}
$B$ = „התוצאות שוות” $=\set{(1,1),\dots,(6,6)}$, 6 תוצאות. נסמן $E$ = „שתי התוצאות זוגיות”. מתוך 6 התוצאות ב-$B$, ב-3 שתיהן זוגיות: $(2,2),(4,4),(6,6)$. במרחב אחיד סופרים בתוך $B$:
\[
\PP(E\mid B)=\frac36=\ans{\frac12}.
\]
\end{solbox}
\item ידוע שבקובייה האדומה יצא 3. מה ההסתברות ששתי התוצאות זוגיות?
\begin{solbox}
אם באדומה יצא 3, היא לא זוגית, ולכן החיתוך ריק וההסתברות היא $\ans{0}$.
\end{solbox}
\end{enumerate}
